\section{Clock 1: depth as time}
\label{sec:clock1}

\subsection{Residual networks as discretized flows}

Equation~\eqref{eq:clock1} is a forward Euler step of an ODE $\dot x = f(x,\ell)$ with unit step. This observation was made early for residual networks \citep{he2015deep} by \citet{e2017proposal}, \citet{haber2017stable} and \citet{lu2017finite}, and it led to continuous-depth models \citep{chen2018neural}, to architectures designed from the stability theory of ODEs and PDEs \citep{haber2017stable,ruthotto2018deep,chang2019antisymmetricrnn}, to invertible second-order variants \citep{sander2021momentum}, and to a view of training as optimal control \citep{e2018meanfield}. Whether a trained residual network actually converges to an ODE as depth grows depends on the initialization and on how the residual branch is scaled \citep{sander2022residual}. In transformers, several stabilization methods amount to a choice of this step size, either by scaling the residual branch at initialization \citep{bachlechner2020rezero,wang2022deepnet} or by moving the layer normalization inside the branch \citep{xiong2020layer,liu2020difficulty}.

\subsection{Attention as an interacting particle system}
\label{sec:particles}

The transformer differs from a plain residual network because the update of token $i$ depends on all tokens. \citet{lu2019improving} read a transformer layer as a splitting scheme for a convection-diffusion equation of interacting particles. \citet{sander2021sinkformers} showed that attention with doubly stochastic normalization is a Wasserstein gradient flow. \citet{geshkovski2023emergence,geshkovski2023mathematical} studied a model in which tokens live on the unit sphere (a stand-in for layer normalization) and follow
\begin{equation}
\dot x_i = P_{x_i}\Big(\frac{1}{Z_i}\sum_{j} e^{\beta\langle Q x_i, K x_j\rangle} V x_j\Big),
\label{eq:particles}
\end{equation}
where $P_x$ projects onto the tangent space and $\beta$ acts as an inverse temperature. They proved that, for broad choices of $Q,K,V$, the tokens cluster as time (depth) grows, in the simplest cases to a single point. Clustering is the particle-system version of an older observation: stacks of pure attention lose rank doubly exponentially with depth \citep{dong2021attention}, and rank collapse at initialization slows training \citep{noci2022signal}. Similar results hold for oversmoothing in attention-based graph networks, which \citet{wu2023demystifying} analyze as nonlinear time-varying dynamical systems.

Later work moved the model closer to real decoders. With causal masking the particle system is no longer a mean-field gradient flow, but clustering results still hold, and the first token, which attends only to itself, plays a special role \citep{karagodin2024clustering}. \citet{wu2024role} showed that the attention mask and layer normalization decide whether collapse happens at all. A mean-field limit with a Vlasov-type ``Transformer PDE'' covers several attention variants \citep{castin2025unified}. Feed-forward layers can be read as a control that steers the particles \citep{maranzatto2026role}. Finite-width transformers with MLP blocks converge to a stochastic interacting particle system, in which noise can synchronize the tokens \citep{agazzi2026stochastic}. For symmetric weights a spectral mode-selection mechanism decides the long-time state \citep{kuehn2026spectral}, and in two dimensions deep linear transformers reduce to Kuramoto-type models with a rich set of behaviors \citep{li2026diverse}. The kinetic theory of causal attention has also been used to explain why models attend poorly to the middle of long contexts \citep{duerinckx2026kinetic,liu2023lost}.

Real models have finite depth, so the long-time limit of \eqref{eq:particles} describes them only partly. \citet{geshkovski2024dynamic} proved that particles stay near a configuration of several clusters for an exponentially long time before they merge, and a stack of finite depth spends its layers inside such a long-lived transient.

\subsection{Signal propagation, order and chaos}
\label{sec:edgechaos}

A second line of work treats a randomly initialized deep network as a random dynamical system and asks whether small input differences grow or shrink with depth. \citet{poole2016exponential} found an order-to-chaos transition controlled by the weight and bias variances; \citet{schoenholz2016deep} showed that the depth over which signals propagate diverges at the transition, so networks are trainable to large depth only near it. The idea goes back to the chaotic phase of random recurrent networks \citep{sompolinsky1988chaos}. Requiring all singular values of the input-output Jacobian to be close to one (dynamical isometry) gives faster training still \citep{pennington2017resurrecting} and allowed 10{,}000-layer convolutional networks to train \citep{xiao2018dynamical}. For transformers, \citet{cowsik2024geometric} track the geometry of $n$ token representations as a discrete-time system of interacting particles and derive conditions on the initialization that make deep transformers trainable.

For trained networks, \citet{storm2023finitetime} computed finite-time Lyapunov exponents \eqref{eq:ftle} of fully connected classifiers and found ridges of large exponents in input space that separate the classes, much like coherent structures in fluid flows. Our first experiment (Section~\ref{sec:exp1}) applies the same quantity to the residual stream of language models across training.

\subsection{Trajectories in trained language models}

Interpretability research gives a complementary, empirical picture of the depth dynamics in trained models. Predictions decoded from intermediate layers improve gradually, which suggests iterative refinement \citep{belrose2023eliciting}. Deleting or swapping middle layers hurts much less than touching the first or last ones, consistent with distinct stages of inference \citep{lad2024remarkable}. The intrinsic dimension of hidden representations rises and then falls across layers \citep{valeriani2023geometry}. Token trajectories through depth cluster along a low-dimensional manifold and can be modeled as a stochastic process with few parameters \citep{sarfati2024lines}. \citet{fernando2025transformer} treat the residual stream with tools from neural population dynamics and report that individual units change smoothly from layer to layer; \citet{fernando2026dynamics} compute full Jacobian eigendecompositions of each layer in production-scale models. A small number of activations are orders of magnitude larger than the rest and nearly independent of the input \citep{sun2024massive}; together with attention sinks \citep{xiao2023efficient}, they can be read as a near fixed point of the depth dynamics. Methods that read features off the layer Jacobian have also been analyzed theoretically \citep{ostrow2026mathematical}. Recurrent-depth or looped models make the iteration explicit by applying the same block repeatedly, and their stability can be studied and scaled as such \citep{prairie2026parcae}.

\subsection{When depth imitates training}

Clock~1 can emulate clock~2. For in-context regression, a linear self-attention layer can implement one step of gradient descent on the in-context examples, so a stack of layers runs an optimizer whose iterations are the layers \citep{oswald2022transformers,akyurek2022learning,garg2022transformers}; trained linear transformers learn preconditioned versions of this algorithm \citep{ahn2023transformers}, and the gradient-flow dynamics that produce them can be solved in closed form \citep{zhang2023trained}. In this regime layer time is optimization time, and the step-size and stability questions of clock~2 carry over to clock~1.

\subsection{Lessons from recurrent networks}

Recurrent networks are the older case in which clocks~1 and~3 coincide: the same map is applied at every time step. Their analysis set many of the patterns used later for transformers. Fixed points and slow points of trained networks explain how they compute \citep{sussillo2013opening}, and sentiment classifiers turn out to integrate evidence along a line attractor \citep{maheswaranathan2019reverse}. On the training side, exploding gradients were linked to bifurcations of the network dynamics \citep{pascanu2012difficulty,ribeiro2019exploding}, and networks trained to reproduce chaotic dynamics must have diverging gradients \citep{mikhaeil2021difficulty}. Lyapunov spectra quantify how much information a recurrent network keeps \citep{engelken2020lyapunov,vogt2020lyapunov}.

\subsection{Explicit dynamical models of sequences}

State-space models write the sequence dynamics in clock~3 explicitly, as a linear dynamical system in a hidden state, with eigenvalues placed to control memory and stability \citep{gu2021efficiently,orvieto2023resurrecting}, or with input-dependent transitions \citep{gu2023mamba}. Linear attention has the same structure: it is a recurrent update of a matrix-valued state \citep{katharopoulos2020transformers,schlag2021linear}. Recurrent networks built from coupled oscillators or constrained by Lipschitz bounds use the same design logic \citep{rusch2020coupled,erichson2020lipschitz}. In these models the attractor and stability questions of this paper are explicit, since the eigenvalues of the state transition are parameters of the architecture.
